\documentclass[aps, reprint, prl, twocolumn, superscriptaddress,amsmath,amssymb,noeprint]{revtex4-2}
% \documentclass{article}
\usepackage{times}
\usepackage{amsmath}
\usepackage{amssymb}
\usepackage{xfrac}
\usepackage{dsfont}
\usepackage{graphicx}
\usepackage{float}
\usepackage{braket}
\usepackage{bbold}
\usepackage[normalem]{ulem}
\usepackage{bm}
\usepackage{bbm}
\usepackage{color}
\usepackage{xcolor}


\usepackage{pdfpages}
\usepackage{pgffor}
\setlength{\parskip}{0pt}
\setlength{\parindent}{1em}

\makeatletter
\AtBeginDocument{\let\LS@rot\@undefined}
\makeatother


\usepackage[colorlinks=true, urlcolor=blue, linkcolor=blue, citecolor=blue]{hyperref}
\usepackage{siunitx}

\newcommand{\ColumbiaEng}{Fu Foundation school of Engineering and Applied Science, Columbia University, New York city, NY}

\frenchspacing

\begin{document}


\title{Quantum Portfolio Optimization for Prediction Markets}

 
\author{Adrian Zhang}
\affiliation{\ColumbiaEng}

\author{Raghav Vishveshwara}
\affiliation{\ColumbiaEng}
\begin{abstract}

We present a portfolio optimization framework that maps prediction market contracts to a Quadratic Unconstrained Binary Optimization (QUBO) problem and hence to an Ising Hamiltonian. Because prediction market contracts resolve to binary outcomes, the mapping is exact: unlike equity portfolios, no binary expansion or granularity parameter is required, and the number of spins equals the number of tradable contracts. Our Hamiltonian minimizes variance risk and maximizes edge while enforcing a budget constraint and a mutual exclusivity penalty. The exclusivity constraint has no analog in equity portfolio optimization, where two assets are never mutually exclusive. Holding several outcomes of one event is a bet on their union, a near-riskless synthetic position when those outcomes cover the event; we show that without the constraint a variance-penalized optimizer fills its limited positions with such combinations instead of directional views. We evaluate the framework on 167 live contracts spanning 32 events from the Kalshi exchange, using exclusivity relations declared by the exchange and verified against settlement rules, and a covariance matrix derived from the martingale property of prediction market prices and estimated from ninety days of price history. Across 100 synthetic forecasts made coherent with the event logic, the penalty-free optimum holds such a combination---for example both candidates in a two-way race---in every draw once risk aversion reaches $\beta = 0.25$, often accepting lower and, at high risk aversion, negative expected edge to do so; the exclusivity penalty removes every such holding. Exhaustive enumeration and a mixed-integer program with hard constraints recover the same ground state, validating the Hamiltonian mapping, while simulated annealing and a linearized-QUBO mixed-integer solver return feasible but suboptimal portfolios. This work establishes prediction markets as a setting in which the Ising formulation is native rather than imposed.
\end{abstract}

\maketitle

\section{Introduction}

Portfolio selection is the central optimization problem of quantitative
finance.  The mean--variance framework introduced by
Markowitz~\cite{markowitz1952}, together with the equilibrium pricing
theory of Sharpe~\cite{sharpe1964} and the analytic characterization of
the efficient frontier by Merton~\cite{merton1972}, expresses the
allocation decision as the optimization of a quadratic objective that
balances expected return against covariance risk.  That quadratic
structure is precisely what makes portfolio selection a natural target
for Ising-model optimization: minimizing a quadratic cost function over
discrete variables is the same computational task as locating the ground
state of a spin Hamiltonian, the problem solved by quantum
annealing~\cite{kadowaki1998,johnson2011} and by variational schemes such
as the Quantum Approximate Optimization Algorithm~\cite{farhi2014}.

The unconstrained mean--variance problem admits a closed-form continuous
solution.  Real trading mandates, however, impose combinatorial
structure: a cap on the number of distinct positions, minimum lot sizes,
and logical restrictions on which assets may be held simultaneously.
Once a cardinality constraint is imposed, the problem becomes a
mixed-integer quadratic program and is
NP-hard~\cite{bienstock1996,bertsimas2009}.  This is the standard
motivation for recasting portfolio selection as a Quadratic Unconstrained
Binary Optimization (QUBO) problem, in which constraints are absorbed
into the objective as quadratic penalty
terms~\cite{lucas2014,glover2019}, and the resulting matrix is mapped to
an Ising Hamiltonian.

\textbf{Combinatorial scaling.}---The practical case for this mapping
rests on how the feasible set grows.  An unconstrained binary portfolio
over $n$ contracts admits $2^n$ configurations.  Imposing a fixed
cardinality $K$ reduces this to $\binom{n}{K}$, which for constant $K$ is
merely polynomial, $\mathcal{O}(n^K)$; exhaustive enumeration therefore
remains viable at small fixed $K$ regardless of $n$.  Realistic mandates,
however, scale the number of held positions with the size of the
investable universe, $K = \lceil \rho n \rceil$ for a roughly constant
participation ratio $\rho$.  The feasible set then grows as
$\binom{n}{\rho n} \sim 2^{n H(\rho)}$, where $H$ is the binary entropy
function---exponential in $n$ (Table~\ref{tab:scaling}).  At a
representative $\rho = 0.05$, exhaustive search is already infeasible by
$n \approx 200$, and the addition of pairwise exclusivity constraints
leaves a constrained combinatorial problem for which no efficient exact
algorithm is known.  Heuristic ground-state search is therefore not a
convenience but a requirement at market scale.

\begin{table}[b]
\caption{\label{tab:scaling}Growth of the portfolio search space.  $K$ is
set by a fixed participation ratio $\rho = 0.05$.  Exhaustive enumeration
becomes infeasible between $n = 100$ and $n = 200$.}
\begin{ruledtabular}
\begin{tabular}{cccc}
$n$ & $2^n$ & $K$ & $\binom{n}{K}$ \\
\colrule
12   & $4.1\times10^{3}$   & 1  & $1.2\times10^{1}$ \\
50   & $1.1\times10^{15}$  & 3  & $2.0\times10^{4}$ \\
100  & $1.3\times10^{30}$  & 5  & $7.5\times10^{7}$ \\
200  & $1.6\times10^{60}$  & 10 & $2.2\times10^{16}$ \\
500  & $3.3\times10^{150}$ & 25 & $1.0\times10^{42}$ \\
1000 & $1.1\times10^{301}$ & 50 & $9.5\times10^{84}$ \\
\end{tabular}
\end{ruledtabular}
\end{table}

\textbf{Quantum and quantum-inspired portfolio optimization.}---A
substantial literature now applies annealing and variational methods to
financial allocation.  Rosenberg \textit{et al.}~\cite{rosenberg2016}
solved the optimal trading trajectory problem on a quantum annealer;
Mugel \textit{et al.}~\cite{mugel2022} performed dynamic portfolio
optimization on real datasets using both quantum processors and
quantum-inspired tensor networks; and broad surveys of the field are
given by Or\'{u}s \textit{et al.}~\cite{orus2019} and Herman \textit{et
al.}~\cite{herman2023}.

A recurring obstacle in this body of work is that equities are
continuously divisible while QUBO variables are binary.  Bridging the two
requires a binary expansion: Lee \textit{et al.}~\cite{lee2024} introduce
a granularity parameter so that several bits encode the fractional weight
of a \emph{single} asset, inflating an $n$-asset problem by that factor.
Beyond the size penalty, the coefficients that enforce the resulting
equality constraints are notoriously sensitive, typically demanding Monte
Carlo estimation or two-stage search to keep the annealer from returning
infeasible solutions while preserving a navigable energy
landscape~\cite{lee2024,verma2022}.  The discretization is, in the words
of the field's own reviews, ``artificial''~\cite{herman2023}: an encoding
artifact rather than a property of the instrument, which makes a genuine
structural advantage over classical continuous solvers difficult to
demonstrate.

\textbf{Prediction markets.}---Prediction markets are exchanges on which
contracts pay a fixed amount contingent on a verifiable event, so that
the traded price is interpretable as a market-aggregated
probability~\cite{wolfers2004,arrow2008}.  Position sizing on such
contracts is classically treated through the Kelly
criterion~\cite{kelly1956,thorp2006}, which maximizes expected log wealth
for an individual bet but does not by itself address discrete selection
under a portfolio-wide capacity limit.  Logical relationships among
contingent claims have been studied in the design of combinatorial
prediction markets~\cite{hanson2003,chen2007}, but there the
\emph{market maker} enforces coherence across a combinatorial outcome
space.  On exchanges such as Kalshi~\cite{kalshi}, contracts within an
event are listed and priced separately and no such coherence is imposed,
so managing positions across mutually exclusive contracts becomes the responsibility
of the \emph{trader's} allocation layer.

To the best of our knowledge, no prior work has applied quantum-inspired
portfolio optimization to prediction markets.  These instruments are
natively binary---an event occurs or it does not, and the investor either
holds the contract or does not---which removes the need for
continuous-to-binary expansion entirely, eliminates the granularity
parameter, and lets the QUBO formulation match the underlying instrument
exactly rather than approximate it.

\textbf{This work.}---We apply this mapping to live market data from the
Kalshi prediction market; we formally designate the resulting algorithm
the Prediction Market Portfolio for Highly Efficient Trading (PMPHET)
framework.  We construct the covariance structure and expected edges and
first evaluate the portfolio without exclusivity constraints.  We show
that, once the objective penalizes risk, the optimizer fills its limited
positions with combinations of mutually exclusive contracts: holding
several sides of one event is a bet on their union, a near-riskless
synthetic position that expresses no view.  We then introduce exclusivity
penalties that remove such holdings, and benchmark our simulated
annealing approach against classical solvers.  The penalty takes the
standard quadratic form used to encode constraints in
QUBOs~\cite{lucas2014}; the constraint itself is, to our knowledge, new
to portfolio optimization: in equity portfolios two assets are never
mutually exclusive, so no formulation in the Markowitz lineage has
required one.  By establishing these
methods classically, we provide a practical foundation for scaling to
larger portfolios and, ultimately, to execution on quantum annealing
hardware.

\textbf{Scope of limitations.}---We emphasize that this framework is
presented as a proof of concept.  We evaluate it on a universe of
$n = 167$ live contracts with a portfolio size of $K = 3$, small enough
that the feasible set can be enumerated exhaustively and every solver
checked against a provably optimal reference.  This isolates the
structural mechanics of the QUBO mapping and the exclusivity penalty.
Establishing where exhaustive and mixed-integer methods cease to be
viable as $K$ grows with $n$, and executing on physical quantum
hardware, are left to subsequent work.

\section{Model}\label{sec:model}\subsection{Problem setup}

Consider $n$ binary prediction contracts with market-implied probabilities
$p_i$. To simulate an external forecasting advantage, we generate model-estimated probabilities $\hat{p}_i$ as a relative perturbation of the market price, $\hat{p}_i = \mathrm{clip}\!\left[p_i(1+u_i),\,0.01,\,0.99\right]$ with $u_i \sim U(-0.10, 0.10)$, so that the perturbation scales with the price rather than displacing near-certain contracts by the same absolute amount as mid-priced ones. We emphasize that this edge is synthetic: the framework treats $\hat{p}_i$ as an input supplied by an external forecasting model, and no claim is made about forecasting performance. The \emph{edge} on
contract~$i$ is $e_i = \hat{p}_i - p_i$.  We seek to select exactly $K$
contracts to maximize total edge subject to portfolio variance risk and
logical consistency constraints.

Independent perturbations ignore the logic of an event.  Within a mutually
exclusive event the forecast probabilities can sum to more than one, and along
a nested ladder a harder threshold can receive a higher probability than an
easier one.  Such a forecast is incoherent and manufactures edge that no
forecaster could hold: in our universe the uncorrected draw assigns probability
$1.06$ to the Federal Reserve cutting rates either zero times or once.  We
therefore project the forecast onto the event logic.  Within each mutually
exclusive event the forecast probabilities are rescaled proportionally
wherever they sum to more than one, and along each nested ladder they are
replaced by their isotonic (monotone) fit in the containment order.  As a
stricter variant we also rescale each mutually exclusive event to the market's
own total, so that no combination spanning an event can carry edge.  Unless
stated otherwise, results use the first projection applied to one fixed draw
of the $u_i$.  The projection removes most of the edge the uncorrected
generator appeared to offer: the valid optimum's edge at $\beta = 0.5$ falls
from $0.115$ to $0.005$.

Each contract $i$ is a Bernoulli random variable with
$\text{Var}(X_i) = p_i(1-p_i)$.  The covariance structure is assembled in
three layers, only the last of which requires estimation.

\textit{(i) Within an event, structure fixes the covariance exactly.}  When
two contracts belong to the same event, their joint distribution is
determined by the event's logic rather than by any correlation
assumption.  For nested thresholds (e.g.\ ``CPI $>0.4\%$'' and ``CPI
$>0.6\%$''), both settling YES requires the harder threshold to be met,
so
\begin{equation}
  \text{Cov}(X_i, X_j) = \min(p_i, p_j) - p_i p_j .
  \label{eq:nested_cov}
\end{equation}
When the event is mutually exclusive, at most one contract settles YES,
$\mathbb{E}[X_iX_j]=0$, and
\begin{equation}
  \text{Cov}(X_i, X_j) = -\,p_i p_j .
  \label{eq:disjoint_cov}
\end{equation}
For $p_i + p_j \le 1$, $-p_ip_j$ is the lower Fr\'{e}chet bound, the most
negative covariance any two binary contracts at those prices can have.
Holding several contracts of one exclusive event therefore offsets risk as
strongly as their prices allow; when the held contracts cover the event, their
combined payoff is identically one and its variance vanishes.  This is what
makes exclusive combinations attractive to a variance-penalized objective
(Sec.~\ref{sec:ablation}).

\textit{(ii) Across events, the covariance is estimated from prices.}  A
prediction market price is the market's running estimate of the outcome,
$P^i_t = \mathbb{E}[X_i \mid \mathcal{F}_t]$ with $P^i_T = X_i$, where the
expectation is taken under the market's pricing measure, and is
therefore a bounded martingale.  Since $P^iP^j - [P^i,P^j]$ is itself a
martingale, where $[\cdot,\cdot]$ denotes the realized quadratic
covariation (which accommodates the jumps that prediction market prices
exhibit, most sharply at settlement), taking expectations at settlement
gives
\begin{equation}
  \text{Cov}(X_i, X_j) \;=\; \mathbb{E}\!\left[\,[P^i, P^j]_T\right],
  \label{eq:covariation}
\end{equation}
so the covariance of the \emph{outcomes} equals the expected quadratic
covariation of the \emph{price paths}.  Equation~\eqref{eq:covariation}
is self-consistent on the diagonal: $\mathbb{E}\big[[P^i]_T\big] =
p_i(1-p_i)$, recovering the Bernoulli variance.  Realized co-movement of
quoted prices is therefore the natural quantity to measure.

Because only part of $[0,T]$ is observed, we estimate realized
\emph{correlation} over the observed window and rescale by the terminal
Bernoulli standard deviation.  This assumes that the correlation of
price increments is stable over each contract's life, and is an
approximation: the contracts in our universe resolve on dates ranging
from September 2026 to 2029, and part of the covariation accrues in the
jump at settlement, which a window preceding resolution does not
observe.  Contracts within one event are driven by a
single underlying quantity and, in a mutually exclusive event, move in
opposite directions; averaging their pairwise correlations against
another event would cancel.  We therefore adopt a one-factor-per-event
model: each event is summarized by a
factor---the first principal component of its contracts' price
increments---with a signed loading $\ell_i$ per contract, and correlation
is estimated between factors.  For $i,j$ in distinct events $e(i),e(j)$,
\begin{equation}
  \Sigma_{ij} = \hat{\rho}_{\,e(i)e(j)}\,\ell_i \ell_j
                \sqrt{p_i(1-p_i)\,p_j(1-p_j)} .
\end{equation}
This reduces the number of quantities to be estimated from $\binom{n}{2}$
contract pairs to $\binom{E}{2}$ event pairs, with $E \ll n$.

\textit{(iii) Shrinkage and positivity.}  With $T$ daily increments each
correlation carries sampling variance $\approx (1-\rho^2)^2/(T-3)$.  With
fewer than a hundred increments per contract and several hundred event
pairs to estimate, much of the observed dispersion is noise, so we
shrink toward a structural prior $F$ fixed in advance (same category
$0.2$, cross category $0$),
\begin{equation}
  \hat{\Sigma} = (1-\delta)\,S + \delta F ,
\end{equation}
with $\delta$ set heuristically as the ratio of the average sampling
variance to the observed variance of the estimates.  The Ledoit--Wolf
intensity, which measures dispersion about the prior rather than about
the sample mean, would assign the data more weight; we retain the more
conservative choice.  Finally, assembling analytic blocks with
estimated cross terms does not guarantee positive semidefiniteness, and a
negative eigenvalue would correspond to a portfolio of negative
variance---a direction the objective actively rewards, since it
\emph{minimizes} $\beta\,\mathbf{x}^\top\Sigma\,\mathbf{x}$.  We therefore
project onto the positive semidefinite cone by clipping negative
eigenvalues and rescaling to restore the Bernoulli diagonal.  The
rescaling perturbs every entry, including the analytic within-event
blocks of (i), which after projection hold only approximately.
Positive semidefiniteness is, moreover, necessary but not sufficient for
binary outcomes: the covariance of two Bernoulli variables must also lie
within the Fr\'{e}chet bounds $\max(0,p_i+p_j-1)-p_ip_j \le
\text{Cov}(X_i,X_j) \le \min(p_i,p_j)-p_ip_j$, which the projected
estimate does not enforce entrywise.

\subsection{QUBO formulation}

The overall objective is to find a binary portfolio vector $\mathbf{x} \in \{0,1\}^n$ that minimizes the following combined cost function:
\begin{equation}
\begin{split}
  \min_{\mathbf{x}} \bigg[ &-\sum_{i} e_i x_i + \beta \sum_{i,j} \Sigma_{ij} x_i x_j \\
  &+ \gamma \left( \sum_i x_i - K \right)^2 + \lambda \sum_{(i,j) \in \mathcal{E}} x_i x_j \bigg].
\end{split}
  \label{eq:objective}
\end{equation}

This cost function maps directly into the standard QUBO form, $\min_{\mathbf{x}} \mathbf{x}^\top Q \, \mathbf{x}$, up to the constant $\gamma K^2$ from expanding the budget penalty, which does not affect the minimizer.  The terms are encoded into the upper-triangular matrix $Q$ as follows:

\textit{(i) Edge maximization.}---Reward selecting high-edge contracts:
\begin{equation}
  Q_{ii} \;\leftarrow\; Q_{ii} - e_i.
\end{equation}

\textit{(ii) Variance penalty.}---Penalize portfolio risk:
\begin{equation}
  Q_{ii} \leftarrow Q_{ii} + \beta\,\Sigma_{ii}, \quad
  Q_{ij} \leftarrow Q_{ij} + 2\beta\,\Sigma_{ij} \;\; (i<j).
\end{equation}

\textit{(iii) Budget constraint.}---Enforce $\sum_i x_i = K$ via a
quadratic penalty $\gamma(\sum_i x_i - K)^2$. Since $x_i^2 = x_i$ for binary variables, expanding this penalty maps the linear components to the diagonal and the interaction components to the off-diagonal:
\begin{equation}
  Q_{ii} \leftarrow Q_{ii} + \gamma(1-2K), \quad
  Q_{ij} \leftarrow Q_{ij} + 2\gamma \;\; (i<j).
\end{equation}

\textit{(iv) Exclusivity constraint.}---For mutually exclusive pairs
$(i,j) \in \mathcal{E}$:
\begin{equation}
  Q_{ij} \leftarrow Q_{ij} + \lambda.
  \label{eq:exclusivity}
\end{equation}
By placing a large penalty $\lambda$ on the off-diagonal, we ensure that the state $x_i = 1$ and $x_j = 1$ becomes energetically highly unfavorable, preventing the solver from holding two outcomes of the same mutually exclusive event.

The resulting $Q$ is mapped to an Ising Hamiltonian
$H = \sum_i h_i \sigma_i^z + \sum_{i<j} J_{ij}\sigma_i^z\sigma_j^z$
via the substitution $x_i = (1+\sigma_i^z)/2$.  The ground state of $H$
corresponds to the optimal portfolio, which we locate by simulated
annealing~\cite{kirkpatrick1983}.
\subsection{Parameters}

We use $\beta = 0.5$ (risk aversion), $\gamma = 2.0$ (budget penalty),
$\lambda = 5.0$ (exclusivity penalty), and $K = 3$ (portfolio size).
The exclusivity penalty $\lambda$ and risk aversion $\beta$ are swept in
Sec.~\ref{sec:ablation}; $\gamma$ is fixed at a
value large enough that every annealing and enumeration solution reported
satisfies the budget constraint.  We fix $K=3$
rather than following the participation-ratio rule of
Table~\ref{tab:scaling} so that the feasible set $\binom{n}{K}$ remains
small enough to enumerate exhaustively, which supplies a provably optimal
reference against which every solver in Sec.~\ref{sec:benchmark} is
measured.

\section{Data}

We construct our universe from open markets on the Kalshi
exchange~\cite{kalshi}, retrieved through its public REST API on
September 8, 2026.  Walking the first $8{,}000$ open events returned by the
events endpoint yields $62{,}945$ contracts carrying a usable quote, grouped
into $7{,}896$ events across nineteen categories.

\textbf{Exclusivity is declared by the exchange.}  Kalshi annotates each event
with a \texttt{mutually\_exclusive} flag, which is set for $2{,}785$ events
covering $18{,}544$ contracts.  We take the set $\mathcal{E}$ from this
flag rather than inferring it from prices.  This matters: a heuristic that
infers exclusivity from ticker strings or from the sum of quoted prices will
misclassify contracts whose direction is not encoded in the ticker---for
instance a ``less than'' threshold listed alongside ``greater than'' thresholds
in the same event---and will assign such a pair a positive covariance when the
truth is Eq.~\eqref{eq:disjoint_cov}.  The flag is not exhaustive, however.
Among our selected events, the market on Venezuela's head of state on
December 31, 2026 is unflagged although at most one candidate can hold the
office; we add it to $\mathcal{E}$ on the basis of its settlement rules.  Its
quoted prices sum to $1.16$, outside the coherence tolerance applied below,
but exclusivity is a property of the settlement logic rather than of the
quotes.  Likewise, five unflagged events whose markets are cumulative
deadlines (e.g.\ ``before January 1, 2027'' and ``before January 20, 2029'')
or price thresholds without a strike in the ticker are classified as nested
from their settlement rules.

\textbf{Selection.}  Contracts are selected event by event rather than
individually, because the structure carrying the covariance information lives at
the event level.  We keep contracts with quotes in $[0.02, 0.98]$ and nonzero
traded volume, retain events with two to fifteen such contracts, and rank
events by traded volume so that price history exists for the estimation of
Eq.~\eqref{eq:covariation}.  For mutually exclusive events we additionally
require the quoted YES prices to sum to unity within $10\%$; roughly $40\%$ of
flagged events fail this no-arbitrage check, typically because untraded
contracts carry stale bid--ask midpoints, and such events carry no usable
covariance structure.  Applying these filters and taking the $200$ most liquid
contracts, then retaining those whose price history spans at least sixty of the
ninety days, leaves $n = 167$ contracts across $32$ events.  Contracts dropped
by these filters are removed from within their events; since
Eqs.~\eqref{eq:nested_cov} and~\eqref{eq:disjoint_cov} hold pairwise, the
analytic covariance of the contracts that remain is unaffected.

\textbf{Structure.}  The resulting universe contains $26$ mutually exclusive
events, generating $|\mathcal{E}| = 511$ exclusivity pairs, and six nested
ladders governed by Eq.~\eqref{eq:nested_cov}; every event therefore has an
analytic within-event form (Table~\ref{tab:contracts}).  The category
composition is given in Table~\ref{tab:categories}.  Of
the $\binom{167}{2} = 13{,}861$ off-diagonal covariance entries, $560$ are fixed
by event structure and the remaining $13{,}301$ are estimated from
$\binom{32}{2} = 496$ event-pair correlations.

\textbf{Price history.}  We retrieve up to ninety daily candles per contract
(June 11 to September 8, 2026) and align them on a common calendar, carrying the
last price forward on days without a close, which gives $T = 89$ price
increments.  Many contracts trade intermittently---the median contract's price
moves on only $49$ of the $89$ days---so the effective sample is smaller than
$T$.  The nominal sampling standard error on a correlation is
$1/\sqrt{T-3} = 0.108$, against an observed dispersion of $0.121$ in
the raw estimates: the data are consistent with only a small amount of genuine
cross-event structure.  Stale prices also bias estimated correlations toward
zero, so this comparison should be read as indicative.  The heuristic
shrinkage weight of Sec.~\ref{sec:model} is $\delta = 0.771$, so the structural prior carries $77\%$ of the final estimate
and the price history $23\%$.  Figure~\ref{fig:correlation} shows the raw and
shrunk correlation matrices on a common scale.  The block structure visible
after shrinkage reflects the same-category prior by construction; the
departures from it---most prominently the pairing of the September Federal
Reserve decision with the 2026 rate-cut-count market---are what the price
history supports.

\begin{figure*}[t]
\includegraphics[width=\textwidth]{code/Ising/exp_200_real/fig_correlation.pdf}
\caption{\label{fig:correlation}Cross-event correlations estimated from $T=89$
daily price increments, shown on a common colour scale.  (a) raw estimates;
(b) after shrinkage toward a same-category structural prior with $\delta=0.771$
chosen from the ratio of sampling noise to observed dispersion.  Events are
ordered by category and the unit diagonal is masked.  The block structure in (b)
is the prior; departures from it are the contribution of the data.}
\end{figure*}

\begin{table}[b]
\caption{\label{tab:contracts}Composition of the $167$-contract Kalshi universe
($32$ events).  The structure class fixes the within-event covariance
analytically, leaving only cross-event terms to be estimated.}
\begin{ruledtabular}
\begin{tabular}{lccc}
Structure & Events & Contracts & Within-event covariance \\
\colrule
Mutually exclusive & 26 & 143 & $-p_i p_j$ \\
Nested & 6 & 24 & $\min(p_i,p_j)-p_i p_j$ \\
\colrule
Total & 32 & 167 & \\
\end{tabular}
\end{ruledtabular}
\end{table}

\begin{table}[b]
\caption{\label{tab:categories}Category composition of the $167$-contract
universe.  Prices are Kalshi quotes on September 8, 2026 (last trade, else
bid--ask midpoint).}
\begin{ruledtabular}
\begin{tabular}{lccc}
Category & Events & Contracts & Price range \\
\colrule
Sports & 11 & 89 & 0.02--0.87 \\
Elections & 9 & 23 & 0.02--0.96 \\
Financials & 2 & 19 & 0.02--0.29 \\
Entertainment & 1 & 10 & 0.02--0.49 \\
Crypto & 2 & 8 & 0.02--0.24 \\
Science and Technology & 2 & 7 & 0.02--0.22 \\
Economics & 2 & 5 & 0.02--0.91 \\
Politics & 2 & 4 & 0.04--0.19 \\
Companies & 1 & 2 & 0.04--0.81 \\
\end{tabular}
\end{ruledtabular}
\end{table}

\section{Results}

\subsection{The Necessity of the Exclusivity Penalty}
\label{sec:ablation}

Holding several contracts from one mutually exclusive event is not a
contradiction in itself: it is a bet that one of them occurs.  Its payoff is
the indicator of their union, whose variance, by
Eq.~\eqref{eq:disjoint_cov}, is $q(1-q)$ with $q = \sum_i p_i$---close to zero
when the contracts cover most of the event's outcomes.  A set of exclusive
contracts whose prices sum to one pays one dollar for a price of one dollar: a
synthetic cash position that expresses no view and, at real bid and ask quotes,
loses the spread.  Under a mandate to hold exactly $K$ positions such
combinations occupy capacity intended for directional views, yet a
variance-penalized objective finds them attractive.  The exclusivity penalty
of Eq.~\eqref{eq:exclusivity} rules them out.

To measure how much this matters we solve the model with the penalty off
($\lambda = 0$) and on ($\lambda \ge 1$) across a range of risk aversion
$\beta$, each case exactly over all $\binom{167}{3}$ portfolios.  Because the
penalty term vanishes whenever no forbidden pair is held, every
$\lambda \ge 1$ defines the same problem; we verified that
$\lambda \in \{1, 2, 5, 10, 20\}$ return the same portfolio.  To separate the
effect from a particular forecast, we repeat the comparison over $100$
independent forecast draws, each projected onto the event logic as described
in Sec.~\ref{sec:model}.

The results are in Table~\ref{tab:ablation}.  At $\beta = 0$ the objective
rewards edge alone and the penalty-free optimum essentially never holds an
exclusive combination.  As soon as risk enters, it does: the share of forecast
draws in which it holds one rises from $6\%$ at $\beta = 0.05$ to $96\%$ at
$\beta = 0.15$ and to $100\%$ from $\beta = 0.25$ onward.  This is the same
interval over which the valid optimum trades edge for variance, from an edge of
$0.092$ at variance $0.60$ ($\beta = 0.1$) to an edge of $0.005$ at variance
$0.059$ ($\beta = 0.2$), after which it settles on three contracts priced at
the $0.02$ selection floor.  The stricter forecast projection and the
uncorrected forecast show the same pattern (caption of
Table~\ref{tab:ablation}).

For the forecast draw used throughout, the penalty-free optimum at the
working value $\beta = 0.5$ buys both candidates in the Michigan Senate race,
priced at $0.63$ and $0.37$.  Its variance is a third of the valid optimum's,
$0.020$ against $0.059$, while its edge is lower, $+0.0018$ against $+0.0051$:
the optimizer gives up return to buy the riskless pair.  At $\beta = 2$ it
holds all three outcomes of the Federal Reserve's September decision, whose
prices sum to $1.02$, reaching a variance of $0.001$ at a \emph{negative}
edge of $-0.020$; across the $100$ coherent draws the penalty-free optimum at
$\beta = 2$ has negative edge in $85\%$ of cases.  It is paying to hold cash.
With the penalty on, no exclusive combination is held in any draw at any
$\beta$.

The effect does not depend on the forecast being coherent.  The uncorrected
generator of Sec.~\ref{sec:model} adds phantom edge on top of it: it assigns
probability $1.06$ to the Federal Reserve cutting rates either zero times or
once, and under that forecast the penalty-free optimum holds exactly that pair.
Projecting the forecast onto the event logic removes the artifact but not the
behaviour.  Neither has a counterpart in equity portfolio optimization, where
two assets are never mutually exclusive and no formulation in the Markowitz
lineage has had cause to forbid a pair.  Both arise specifically because
prediction market contracts are event-contingent claims.

\begin{table*}[t]
\caption{\label{tab:ablation}Exclusivity ablation across risk aversion on
$167$ live Kalshi contracts with $K = 3$, each entry solved exactly over all
$\binom{167}{3}$ portfolios.  Penalty on: $\lambda \ge 1$ (all of
$\lambda \in \{1,2,5,10,20\}$ give the same portfolio).  Penalty off:
$\lambda = 0$.  Edge and variance are for the forecast draw used throughout,
projected onto the event logic (Sec.~\ref{sec:model}).  The last column is the
share of $100$ independent projected forecast draws whose penalty-free optimum
holds a mutually exclusive combination; with the penalty on the share is zero
at every $\beta$.  Under the stricter projection that matches the market's
event totals the shares are $0, 2, 12, 86, 99, 100, 100, 100, 100\%$; under the
uncorrected independent-noise forecast they are
$3, 23, 52, 73, 84, 91, 100, 100, 100\%$.}
\begin{ruledtabular}
\begin{tabular}{ccccccc}
 & \multicolumn{2}{c}{Penalty on} & \multicolumn{3}{c}{Penalty off} & \\
$\beta$ & Edge & Variance & Edge & Variance & Exclusive holding & Share of draws \\
\colrule
0    & $+0.0915$ & 0.6239 & $+0.0915$ & 0.6239 & ---                       & $0\%$ \\
0.05 & $+0.0915$ & 0.6032 & $+0.0915$ & 0.6032 & ---                       & $6\%$ \\
0.1  & $+0.0915$ & 0.6032 & $+0.0915$ & 0.6032 & ---                       & $60\%$ \\
0.15 & $+0.0731$ & 0.4604 & $+0.0434$ & 0.2491 & Michigan Senate, D and R  & $96\%$ \\
0.2  & $+0.0051$ & 0.0586 & $+0.0018$ & 0.0196 & Michigan Senate, D and R  & $99\%$ \\
0.25 & $+0.0051$ & 0.0586 & $+0.0018$ & 0.0196 & Michigan Senate, D and R  & $100\%$ \\
\textbf{0.5} & $+0.0051$ & 0.0586 & $+0.0018$ & \textbf{0.0196} & Michigan Senate, D and R & $100\%$ \\
1    & $+0.0051$ & 0.0586 & $+0.0018$ & 0.0196 & Michigan Senate, D and R  & $100\%$ \\
2    & $+0.0033$ & 0.0571 & $-0.0200$ & 0.0011 & Fed September, all three  & $100\%$ \\
\end{tabular}
\end{ruledtabular}
\end{table*}

\subsection{Validating the Hamiltonian Mapping}
\label{sec:benchmark}

Because $K=3$, the feasible set contains $\binom{167}{3} = 762{,}355$
portfolios, which can be enumerated exhaustively in about five seconds.  This
supplies a provably optimal reference, and we use it to check how independent
solution methods applied to the same Hamiltonian compare with its ground state
(Table~\ref{tab:benchmark}).  Energies are reported with the constant
$\gamma K^2$ restored, so that they equal the objective of
Eq.~\eqref{eq:objective}; for a feasible portfolio this is
$-(\text{edge}) + \beta\,(\text{variance})$.

Exhaustive enumeration and the mixed-integer program with hard
constraints both return energy $+0.024132$ on the same three contracts.
Simulated annealing returns a feasible portfolio at $+0.026021$, and the
linearized QUBO handed to a branch-and-cut solver reaches $+0.028765$ before
exhausting a 300\,s budget.  The exact agreement of two methods with no
algorithmic structure in common is evidence that the QUBO of
Eq.~\eqref{eq:objective} encodes the intended problem; the gaps of the other
two are properties of those solvers, not of the mapping.

The optimal energy is positive: at $\beta = 0.5$ even the best valid
three-contract portfolio has a negative risk-adjusted value, its edge of
$0.0051$ falling short of $\beta$ times its variance of $0.0586$.  An
unconstrained investor would hold nothing.  The budget term encodes a mandate
to hold exactly $K$ positions, and the optimum is the least costly way to
satisfy it; this is the regime in which the synthetic cash positions of
Sec.~\ref{sec:ablation} are most tempting.

Two observations qualify any performance reading of this table.  First,
simulated annealing does \emph{not} recover the optimum at this scale:
it is slower than exhaustive enumeration, and its gap of $0.0019$ is about
$8\%$ of the optimal objective.  It does, however, respect every
constraint---budget and exclusivity---in every run.  A plausible contributor
is the scale of the penalty terms: $\gamma$ and $\lambda$ exceed the objective
differences between good feasible portfolios by roughly three orders of
magnitude, so a temperature schedule set by the largest couplings resolves
feasibility well but ranks feasible portfolios poorly.  Second, continuous
relaxation followed by rounding fails
outright, returning an empty portfolio that satisfies neither the budget
constraint nor any notion of optimality.  The natively binary structure of
prediction contracts is not amenable to relaxation, which is precisely the
argument for a formulation in which the binary variables are the instruments
themselves.

\begin{table}[b]
\caption{\label{tab:benchmark}Solver comparison on $167$ live Kalshi contracts
with $K=3$.  Energy is the objective of Eq.~\eqref{eq:objective}, including
the constant $\gamma K^2$; lower is better.  Gap is $E - E^\ast$ against
exhaustive enumeration over the $762{,}355$ feasible portfolios.  Both
mixed-integer encodings require $13{,}861$ auxiliary binary variables; the
linearized QUBO did not converge within its time limit.  The relaxation selects
no contracts, so its energy is the budget penalty $\gamma K^2 = 18$.}
\begin{ruledtabular}
\begin{tabular}{lccccr}
Solver & Energy & Gap & Budget & Viol. & Time (s) \\
\colrule
Exhaustive $\binom{n}{K}$ & $+0.024132$ & ---        & $\checkmark$ & 0 & 5.01 \\
PuLP+CBC (native)         & $+0.024132$ & $0$        & $\checkmark$ & 0 & 217.82 \\
SA (Ising, 5000 reads)    & $+0.026021$ & $0.0019$   & $\checkmark$ & 0 & 35.67 \\
PuLP+CBC (QUBO)           & $+0.028765$ & $0.0046$   & $\checkmark$ & 0 & 301.37 \\
SciPy relax+round         & $+18.0$     & ---        & $\times$     & 0 & 5.36 \\
\end{tabular}
\end{ruledtabular}
\end{table}

\subsection{Simulated Annealing Across Seeds}
\label{sec:robustness}

Across twenty independent annealing runs with distinct random seeds, every run
satisfied the budget constraint and none held an exclusive pair.  The runs
returned nineteen \emph{distinct} portfolios---mean energy $+0.0259$, standard
deviation $0.0006$, best $+0.0245$---and none reached the exact optimum of
$+0.0241$; every run finished between $0.0004$ and $0.0031$ above it.  At this
scale, simulated annealing with $5000$ reads reliably finds feasible portfolios
close to the optimum, but not the optimum itself.

\section{Conclusion}

We have presented a portfolio optimization framework in which the mapping to an
Ising Hamiltonian is exact rather than approximate.  Equity formulations must
discretize a continuous weight vector, spending several spins per asset on a
granularity parameter that is an artifact of the encoding; prediction market
contracts are already binary, so the number of spins equals the number of
tradable instruments and the decision space is represented without loss.

The framework's distinctive constraint is exclusivity.  Holding several
outcomes of one event is a bet on their union, and when those outcomes cover
the event the position is synthetic cash.  We showed that once risk enters the
objective, an optimizer without the constraint fills its positions with such
combinations---in every one of $100$ forecast draws for $\beta \ge 0.25$---often
accepting lower and, at high risk aversion, negative edge to do so.  In our live
universe it bought both candidates in the Michigan Senate race and, at
$\beta = 2$, all three outcomes of the Federal Reserve's September decision.
The exclusivity penalty removes every such holding.  We also found that a
forecast generated without regard to event logic manufactures edge on exactly
these combinations; projecting it onto the event structure removes the
artifact but not the behaviour.  Since equity assets are never mutually
exclusive, no formulation in the Markowitz lineage has needed such a
constraint; it is specific to event-contingent claims.

We also replaced the hand-set correlation matrix usual in this setting with an
estimate derived from the martingale property of prediction market prices,
Eq.~\eqref{eq:covariation}, and measured from price history.  On ninety
days of data the sampling error is comparable to the dispersion of the estimates
themselves, and the heuristic shrinkage weight is $\delta = 0.771$:
the structural prior does most of the work, and we report this rather than
overstate what the price history supports.

Exhaustive enumeration and a mixed-integer program with hard constraints,
methods with no algorithmic structure in common, recover the same ground state
of the resulting Hamiltonian, which validates the mapping.  We caution against
reading Table~\ref{tab:benchmark} as a performance claim: at $n=167$ with $K=3$,
exhaustive enumeration is the fastest method, and simulated annealing, while
respecting every constraint, did not recover the exact optimum in any of twenty
independent runs.  Rescaling the penalty terms relative to the objective is a
natural first remedy.  The case for the Ising formulation rests not
on solver speed at this scale but on the exactness of the encoding and on the
growth of the feasible set (Table~\ref{tab:scaling}), which removes exhaustive
search as an option once the number of held positions scales with the size of
the investable universe.  Establishing where that crossover occurs empirically,
and executing the same Hamiltonian on quantum annealing hardware, are the
natural next steps.

\section*{Acknowledgments}
The authors fondly dedicate the PMPHET acronym to their cats: Puma, Poppy, Hershey, Elliot, and Tiger, whose quiet support was foundational to the completion of this work.

\begin{thebibliography}{30}

\bibitem{markowitz1952}
H.~Markowitz, \emph{J.\ Finance} \textbf{7}, 77 (1952).

\bibitem{sharpe1964}
W.~F.~Sharpe, \emph{J.\ Finance} \textbf{19}, 425 (1964).

\bibitem{merton1972}
R.~C.~Merton, \emph{J.\ Financial Quant.\ Anal.} \textbf{7}, 1851 (1972).

\bibitem{kadowaki1998}
T.~Kadowaki and H.~Nishimori, \emph{Phys.\ Rev.\ E} \textbf{58}, 5355 (1998).

\bibitem{johnson2011}
M.~W.~Johnson \emph{et al.}, \emph{Nature} \textbf{473}, 194 (2011).

\bibitem{farhi2014}
E.~Farhi, J.~Goldstone, and S.~Gutmann, \emph{A Quantum Approximate
Optimization Algorithm}, arXiv:1411.4028 (2014).

\bibitem{bienstock1996}
D.~Bienstock, \emph{Math.\ Program.} \textbf{74}, 121 (1996).

\bibitem{bertsimas2009}
D.~Bertsimas and R.~Shioda, \emph{Comput.\ Optim.\ Appl.} \textbf{43}, 1 (2009).

\bibitem{lucas2014}
A.~Lucas, \emph{Front.\ Phys.} \textbf{2}, 5 (2014).

\bibitem{glover2019}
F.~Glover, G.~Kochenberger, and Y.~Du, \emph{4OR} \textbf{17}, 335 (2019).


\bibitem{rosenberg2016}
G.~Rosenberg \emph{et al.}, \emph{IEEE J.\ Sel.\ Top.\ Signal Process.}
\textbf{10}, 1053 (2016).

\bibitem{mugel2022}
S.~Mugel \emph{et al.}, \emph{Phys.\ Rev.\ Research} \textbf{4}, 013006 (2022).

\bibitem{orus2019}
R.~Or\'{u}s, S.~Mugel, and E.~Lizaso, \emph{Rev.\ Phys.} \textbf{4}, 100028
(2019).

\bibitem{herman2023}
D.~Herman \emph{et al.}, \emph{Nat.\ Rev.\ Phys.} \textbf{5}, 450 (2023).

\bibitem{lee2024}
Lee \textit{et al.}, \emph{Quantum-Inspired Portfolio Optimization In The QUBO
Framework}, arXiv:2410.05932 (2024).

\bibitem{verma2022}
A.~Verma and M.~Lewis, \emph{Discrete Optim.} \textbf{44}, 100594 (2022).

\bibitem{wolfers2004}
J.~Wolfers and E.~Zitzewitz, \emph{J.\ Econ.\ Perspect.} \textbf{18}, 107
(2004).

\bibitem{arrow2008}
K.~J.~Arrow \emph{et al.}, \emph{Science} \textbf{320}, 877 (2008).

\bibitem{kelly1956}
J.~L.~Kelly~Jr., \emph{Bell Syst.\ Tech.\ J.} \textbf{35}, 917 (1956).

\bibitem{thorp2006}
E.~O.~Thorp, in \emph{Handbook of Asset and Liability Management},
Vol.~1 (Elsevier, Amsterdam, 2006), p.~385.

\bibitem{hanson2003}
R.~Hanson, \emph{Inf.\ Syst.\ Front.} \textbf{5}, 107 (2003).

\bibitem{chen2007}
Y.~Chen and D.~M.~Pennock, in \emph{Proceedings of the 23rd Conference on
Uncertainty in Artificial Intelligence (UAI)} (2007), p.~49.

\bibitem{kalshi}
Kalshi Exchange, \url{https://kalshi.com} (2026).


\bibitem{kirkpatrick1983}
S.~Kirkpatrick, C.~D.~Gelatt, and M.~P.~Vecchi, \emph{Science} \textbf{220},
671 (1983).

\end{thebibliography}
\end{document}
